
% Default to the notebook output style

    


% Inherit from the specified cell style.




    
\documentclass[11pt]{article}

    
    
    \usepackage[T1]{fontenc}
    % Nicer default font (+ math font) than Computer Modern for most use cases
    \usepackage{mathpazo}

    % Basic figure setup, for now with no caption control since it's done
    % automatically by Pandoc (which extracts ![](path) syntax from Markdown).
    \usepackage{graphicx}
    % We will generate all images so they have a width \maxwidth. This means
    % that they will get their normal width if they fit onto the page, but
    % are scaled down if they would overflow the margins.
    \makeatletter
    \def\maxwidth{\ifdim\Gin@nat@width>\linewidth\linewidth
    \else\Gin@nat@width\fi}
    \makeatother
    \let\Oldincludegraphics\includegraphics
    % Set max figure width to be 80% of text width, for now hardcoded.
    \renewcommand{\includegraphics}[1]{\Oldincludegraphics[width=.8\maxwidth]{#1}}
    % Ensure that by default, figures have no caption (until we provide a
    % proper Figure object with a Caption API and a way to capture that
    % in the conversion process - todo).
    \usepackage{caption}
    \DeclareCaptionLabelFormat{nolabel}{}
    \captionsetup{labelformat=nolabel}

    \usepackage{adjustbox} % Used to constrain images to a maximum size 
    \usepackage{xcolor} % Allow colors to be defined
    \usepackage{enumerate} % Needed for markdown enumerations to work
    \usepackage{geometry} % Used to adjust the document margins
    \usepackage{amsmath} % Equations
    \usepackage{amssymb} % Equations
    \usepackage{textcomp} % defines textquotesingle
    % Hack from http://tex.stackexchange.com/a/47451/13684:
    \AtBeginDocument{%
        \def\PYZsq{\textquotesingle}% Upright quotes in Pygmentized code
    }
    \usepackage{upquote} % Upright quotes for verbatim code
    \usepackage{eurosym} % defines \euro
    \usepackage[mathletters]{ucs} % Extended unicode (utf-8) support
    \usepackage[utf8x]{inputenc} % Allow utf-8 characters in the tex document
    \usepackage{fancyvrb} % verbatim replacement that allows latex
    \usepackage{grffile} % extends the file name processing of package graphics 
                         % to support a larger range 
    % The hyperref package gives us a pdf with properly built
    % internal navigation ('pdf bookmarks' for the table of contents,
    % internal cross-reference links, web links for URLs, etc.)
    \usepackage{hyperref}
    \usepackage{longtable} % longtable support required by pandoc >1.10
    \usepackage{booktabs}  % table support for pandoc > 1.12.2
    \usepackage[inline]{enumitem} % IRkernel/repr support (it uses the enumerate* environment)
    \usepackage[normalem]{ulem} % ulem is needed to support strikethroughs (\sout)
                                % normalem makes italics be italics, not underlines
    

    
    
    % Colors for the hyperref package
    \definecolor{urlcolor}{rgb}{0,.145,.698}
    \definecolor{linkcolor}{rgb}{.71,0.21,0.01}
    \definecolor{citecolor}{rgb}{.12,.54,.11}

    % ANSI colors
    \definecolor{ansi-black}{HTML}{3E424D}
    \definecolor{ansi-black-intense}{HTML}{282C36}
    \definecolor{ansi-red}{HTML}{E75C58}
    \definecolor{ansi-red-intense}{HTML}{B22B31}
    \definecolor{ansi-green}{HTML}{00A250}
    \definecolor{ansi-green-intense}{HTML}{007427}
    \definecolor{ansi-yellow}{HTML}{DDB62B}
    \definecolor{ansi-yellow-intense}{HTML}{B27D12}
    \definecolor{ansi-blue}{HTML}{208FFB}
    \definecolor{ansi-blue-intense}{HTML}{0065CA}
    \definecolor{ansi-magenta}{HTML}{D160C4}
    \definecolor{ansi-magenta-intense}{HTML}{A03196}
    \definecolor{ansi-cyan}{HTML}{60C6C8}
    \definecolor{ansi-cyan-intense}{HTML}{258F8F}
    \definecolor{ansi-white}{HTML}{C5C1B4}
    \definecolor{ansi-white-intense}{HTML}{A1A6B2}

    % commands and environments needed by pandoc snippets
    % extracted from the output of `pandoc -s`
    \providecommand{\tightlist}{%
      \setlength{\itemsep}{0pt}\setlength{\parskip}{0pt}}
    \DefineVerbatimEnvironment{Highlighting}{Verbatim}{commandchars=\\\{\}}
    % Add ',fontsize=\small' for more characters per line
    \newenvironment{Shaded}{}{}
    \newcommand{\KeywordTok}[1]{\textcolor[rgb]{0.00,0.44,0.13}{\textbf{{#1}}}}
    \newcommand{\DataTypeTok}[1]{\textcolor[rgb]{0.56,0.13,0.00}{{#1}}}
    \newcommand{\DecValTok}[1]{\textcolor[rgb]{0.25,0.63,0.44}{{#1}}}
    \newcommand{\BaseNTok}[1]{\textcolor[rgb]{0.25,0.63,0.44}{{#1}}}
    \newcommand{\FloatTok}[1]{\textcolor[rgb]{0.25,0.63,0.44}{{#1}}}
    \newcommand{\CharTok}[1]{\textcolor[rgb]{0.25,0.44,0.63}{{#1}}}
    \newcommand{\StringTok}[1]{\textcolor[rgb]{0.25,0.44,0.63}{{#1}}}
    \newcommand{\CommentTok}[1]{\textcolor[rgb]{0.38,0.63,0.69}{\textit{{#1}}}}
    \newcommand{\OtherTok}[1]{\textcolor[rgb]{0.00,0.44,0.13}{{#1}}}
    \newcommand{\AlertTok}[1]{\textcolor[rgb]{1.00,0.00,0.00}{\textbf{{#1}}}}
    \newcommand{\FunctionTok}[1]{\textcolor[rgb]{0.02,0.16,0.49}{{#1}}}
    \newcommand{\RegionMarkerTok}[1]{{#1}}
    \newcommand{\ErrorTok}[1]{\textcolor[rgb]{1.00,0.00,0.00}{\textbf{{#1}}}}
    \newcommand{\NormalTok}[1]{{#1}}
    
    % Additional commands for more recent versions of Pandoc
    \newcommand{\ConstantTok}[1]{\textcolor[rgb]{0.53,0.00,0.00}{{#1}}}
    \newcommand{\SpecialCharTok}[1]{\textcolor[rgb]{0.25,0.44,0.63}{{#1}}}
    \newcommand{\VerbatimStringTok}[1]{\textcolor[rgb]{0.25,0.44,0.63}{{#1}}}
    \newcommand{\SpecialStringTok}[1]{\textcolor[rgb]{0.73,0.40,0.53}{{#1}}}
    \newcommand{\ImportTok}[1]{{#1}}
    \newcommand{\DocumentationTok}[1]{\textcolor[rgb]{0.73,0.13,0.13}{\textit{{#1}}}}
    \newcommand{\AnnotationTok}[1]{\textcolor[rgb]{0.38,0.63,0.69}{\textbf{\textit{{#1}}}}}
    \newcommand{\CommentVarTok}[1]{\textcolor[rgb]{0.38,0.63,0.69}{\textbf{\textit{{#1}}}}}
    \newcommand{\VariableTok}[1]{\textcolor[rgb]{0.10,0.09,0.49}{{#1}}}
    \newcommand{\ControlFlowTok}[1]{\textcolor[rgb]{0.00,0.44,0.13}{\textbf{{#1}}}}
    \newcommand{\OperatorTok}[1]{\textcolor[rgb]{0.40,0.40,0.40}{{#1}}}
    \newcommand{\BuiltInTok}[1]{{#1}}
    \newcommand{\ExtensionTok}[1]{{#1}}
    \newcommand{\PreprocessorTok}[1]{\textcolor[rgb]{0.74,0.48,0.00}{{#1}}}
    \newcommand{\AttributeTok}[1]{\textcolor[rgb]{0.49,0.56,0.16}{{#1}}}
    \newcommand{\InformationTok}[1]{\textcolor[rgb]{0.38,0.63,0.69}{\textbf{\textit{{#1}}}}}
    \newcommand{\WarningTok}[1]{\textcolor[rgb]{0.38,0.63,0.69}{\textbf{\textit{{#1}}}}}
    
    
    % Define a nice break command that doesn't care if a line doesn't already
    % exist.
    \def\br{\hspace*{\fill} \\* }
    % Math Jax compatability definitions
    \def\gt{>}
    \def\lt{<}
    % Document parameters
    \title{00\_probabilities}
    
    
    

    % Pygments definitions
    
\makeatletter
\def\PY@reset{\let\PY@it=\relax \let\PY@bf=\relax%
    \let\PY@ul=\relax \let\PY@tc=\relax%
    \let\PY@bc=\relax \let\PY@ff=\relax}
\def\PY@tok#1{\csname PY@tok@#1\endcsname}
\def\PY@toks#1+{\ifx\relax#1\empty\else%
    \PY@tok{#1}\expandafter\PY@toks\fi}
\def\PY@do#1{\PY@bc{\PY@tc{\PY@ul{%
    \PY@it{\PY@bf{\PY@ff{#1}}}}}}}
\def\PY#1#2{\PY@reset\PY@toks#1+\relax+\PY@do{#2}}

\expandafter\def\csname PY@tok@w\endcsname{\def\PY@tc##1{\textcolor[rgb]{0.73,0.73,0.73}{##1}}}
\expandafter\def\csname PY@tok@c\endcsname{\let\PY@it=\textit\def\PY@tc##1{\textcolor[rgb]{0.25,0.50,0.50}{##1}}}
\expandafter\def\csname PY@tok@cp\endcsname{\def\PY@tc##1{\textcolor[rgb]{0.74,0.48,0.00}{##1}}}
\expandafter\def\csname PY@tok@k\endcsname{\let\PY@bf=\textbf\def\PY@tc##1{\textcolor[rgb]{0.00,0.50,0.00}{##1}}}
\expandafter\def\csname PY@tok@kp\endcsname{\def\PY@tc##1{\textcolor[rgb]{0.00,0.50,0.00}{##1}}}
\expandafter\def\csname PY@tok@kt\endcsname{\def\PY@tc##1{\textcolor[rgb]{0.69,0.00,0.25}{##1}}}
\expandafter\def\csname PY@tok@o\endcsname{\def\PY@tc##1{\textcolor[rgb]{0.40,0.40,0.40}{##1}}}
\expandafter\def\csname PY@tok@ow\endcsname{\let\PY@bf=\textbf\def\PY@tc##1{\textcolor[rgb]{0.67,0.13,1.00}{##1}}}
\expandafter\def\csname PY@tok@nb\endcsname{\def\PY@tc##1{\textcolor[rgb]{0.00,0.50,0.00}{##1}}}
\expandafter\def\csname PY@tok@nf\endcsname{\def\PY@tc##1{\textcolor[rgb]{0.00,0.00,1.00}{##1}}}
\expandafter\def\csname PY@tok@nc\endcsname{\let\PY@bf=\textbf\def\PY@tc##1{\textcolor[rgb]{0.00,0.00,1.00}{##1}}}
\expandafter\def\csname PY@tok@nn\endcsname{\let\PY@bf=\textbf\def\PY@tc##1{\textcolor[rgb]{0.00,0.00,1.00}{##1}}}
\expandafter\def\csname PY@tok@ne\endcsname{\let\PY@bf=\textbf\def\PY@tc##1{\textcolor[rgb]{0.82,0.25,0.23}{##1}}}
\expandafter\def\csname PY@tok@nv\endcsname{\def\PY@tc##1{\textcolor[rgb]{0.10,0.09,0.49}{##1}}}
\expandafter\def\csname PY@tok@no\endcsname{\def\PY@tc##1{\textcolor[rgb]{0.53,0.00,0.00}{##1}}}
\expandafter\def\csname PY@tok@nl\endcsname{\def\PY@tc##1{\textcolor[rgb]{0.63,0.63,0.00}{##1}}}
\expandafter\def\csname PY@tok@ni\endcsname{\let\PY@bf=\textbf\def\PY@tc##1{\textcolor[rgb]{0.60,0.60,0.60}{##1}}}
\expandafter\def\csname PY@tok@na\endcsname{\def\PY@tc##1{\textcolor[rgb]{0.49,0.56,0.16}{##1}}}
\expandafter\def\csname PY@tok@nt\endcsname{\let\PY@bf=\textbf\def\PY@tc##1{\textcolor[rgb]{0.00,0.50,0.00}{##1}}}
\expandafter\def\csname PY@tok@nd\endcsname{\def\PY@tc##1{\textcolor[rgb]{0.67,0.13,1.00}{##1}}}
\expandafter\def\csname PY@tok@s\endcsname{\def\PY@tc##1{\textcolor[rgb]{0.73,0.13,0.13}{##1}}}
\expandafter\def\csname PY@tok@sd\endcsname{\let\PY@it=\textit\def\PY@tc##1{\textcolor[rgb]{0.73,0.13,0.13}{##1}}}
\expandafter\def\csname PY@tok@si\endcsname{\let\PY@bf=\textbf\def\PY@tc##1{\textcolor[rgb]{0.73,0.40,0.53}{##1}}}
\expandafter\def\csname PY@tok@se\endcsname{\let\PY@bf=\textbf\def\PY@tc##1{\textcolor[rgb]{0.73,0.40,0.13}{##1}}}
\expandafter\def\csname PY@tok@sr\endcsname{\def\PY@tc##1{\textcolor[rgb]{0.73,0.40,0.53}{##1}}}
\expandafter\def\csname PY@tok@ss\endcsname{\def\PY@tc##1{\textcolor[rgb]{0.10,0.09,0.49}{##1}}}
\expandafter\def\csname PY@tok@sx\endcsname{\def\PY@tc##1{\textcolor[rgb]{0.00,0.50,0.00}{##1}}}
\expandafter\def\csname PY@tok@m\endcsname{\def\PY@tc##1{\textcolor[rgb]{0.40,0.40,0.40}{##1}}}
\expandafter\def\csname PY@tok@gh\endcsname{\let\PY@bf=\textbf\def\PY@tc##1{\textcolor[rgb]{0.00,0.00,0.50}{##1}}}
\expandafter\def\csname PY@tok@gu\endcsname{\let\PY@bf=\textbf\def\PY@tc##1{\textcolor[rgb]{0.50,0.00,0.50}{##1}}}
\expandafter\def\csname PY@tok@gd\endcsname{\def\PY@tc##1{\textcolor[rgb]{0.63,0.00,0.00}{##1}}}
\expandafter\def\csname PY@tok@gi\endcsname{\def\PY@tc##1{\textcolor[rgb]{0.00,0.63,0.00}{##1}}}
\expandafter\def\csname PY@tok@gr\endcsname{\def\PY@tc##1{\textcolor[rgb]{1.00,0.00,0.00}{##1}}}
\expandafter\def\csname PY@tok@ge\endcsname{\let\PY@it=\textit}
\expandafter\def\csname PY@tok@gs\endcsname{\let\PY@bf=\textbf}
\expandafter\def\csname PY@tok@gp\endcsname{\let\PY@bf=\textbf\def\PY@tc##1{\textcolor[rgb]{0.00,0.00,0.50}{##1}}}
\expandafter\def\csname PY@tok@go\endcsname{\def\PY@tc##1{\textcolor[rgb]{0.53,0.53,0.53}{##1}}}
\expandafter\def\csname PY@tok@gt\endcsname{\def\PY@tc##1{\textcolor[rgb]{0.00,0.27,0.87}{##1}}}
\expandafter\def\csname PY@tok@err\endcsname{\def\PY@bc##1{\setlength{\fboxsep}{0pt}\fcolorbox[rgb]{1.00,0.00,0.00}{1,1,1}{\strut ##1}}}
\expandafter\def\csname PY@tok@kc\endcsname{\let\PY@bf=\textbf\def\PY@tc##1{\textcolor[rgb]{0.00,0.50,0.00}{##1}}}
\expandafter\def\csname PY@tok@kd\endcsname{\let\PY@bf=\textbf\def\PY@tc##1{\textcolor[rgb]{0.00,0.50,0.00}{##1}}}
\expandafter\def\csname PY@tok@kn\endcsname{\let\PY@bf=\textbf\def\PY@tc##1{\textcolor[rgb]{0.00,0.50,0.00}{##1}}}
\expandafter\def\csname PY@tok@kr\endcsname{\let\PY@bf=\textbf\def\PY@tc##1{\textcolor[rgb]{0.00,0.50,0.00}{##1}}}
\expandafter\def\csname PY@tok@bp\endcsname{\def\PY@tc##1{\textcolor[rgb]{0.00,0.50,0.00}{##1}}}
\expandafter\def\csname PY@tok@fm\endcsname{\def\PY@tc##1{\textcolor[rgb]{0.00,0.00,1.00}{##1}}}
\expandafter\def\csname PY@tok@vc\endcsname{\def\PY@tc##1{\textcolor[rgb]{0.10,0.09,0.49}{##1}}}
\expandafter\def\csname PY@tok@vg\endcsname{\def\PY@tc##1{\textcolor[rgb]{0.10,0.09,0.49}{##1}}}
\expandafter\def\csname PY@tok@vi\endcsname{\def\PY@tc##1{\textcolor[rgb]{0.10,0.09,0.49}{##1}}}
\expandafter\def\csname PY@tok@vm\endcsname{\def\PY@tc##1{\textcolor[rgb]{0.10,0.09,0.49}{##1}}}
\expandafter\def\csname PY@tok@sa\endcsname{\def\PY@tc##1{\textcolor[rgb]{0.73,0.13,0.13}{##1}}}
\expandafter\def\csname PY@tok@sb\endcsname{\def\PY@tc##1{\textcolor[rgb]{0.73,0.13,0.13}{##1}}}
\expandafter\def\csname PY@tok@sc\endcsname{\def\PY@tc##1{\textcolor[rgb]{0.73,0.13,0.13}{##1}}}
\expandafter\def\csname PY@tok@dl\endcsname{\def\PY@tc##1{\textcolor[rgb]{0.73,0.13,0.13}{##1}}}
\expandafter\def\csname PY@tok@s2\endcsname{\def\PY@tc##1{\textcolor[rgb]{0.73,0.13,0.13}{##1}}}
\expandafter\def\csname PY@tok@sh\endcsname{\def\PY@tc##1{\textcolor[rgb]{0.73,0.13,0.13}{##1}}}
\expandafter\def\csname PY@tok@s1\endcsname{\def\PY@tc##1{\textcolor[rgb]{0.73,0.13,0.13}{##1}}}
\expandafter\def\csname PY@tok@mb\endcsname{\def\PY@tc##1{\textcolor[rgb]{0.40,0.40,0.40}{##1}}}
\expandafter\def\csname PY@tok@mf\endcsname{\def\PY@tc##1{\textcolor[rgb]{0.40,0.40,0.40}{##1}}}
\expandafter\def\csname PY@tok@mh\endcsname{\def\PY@tc##1{\textcolor[rgb]{0.40,0.40,0.40}{##1}}}
\expandafter\def\csname PY@tok@mi\endcsname{\def\PY@tc##1{\textcolor[rgb]{0.40,0.40,0.40}{##1}}}
\expandafter\def\csname PY@tok@il\endcsname{\def\PY@tc##1{\textcolor[rgb]{0.40,0.40,0.40}{##1}}}
\expandafter\def\csname PY@tok@mo\endcsname{\def\PY@tc##1{\textcolor[rgb]{0.40,0.40,0.40}{##1}}}
\expandafter\def\csname PY@tok@ch\endcsname{\let\PY@it=\textit\def\PY@tc##1{\textcolor[rgb]{0.25,0.50,0.50}{##1}}}
\expandafter\def\csname PY@tok@cm\endcsname{\let\PY@it=\textit\def\PY@tc##1{\textcolor[rgb]{0.25,0.50,0.50}{##1}}}
\expandafter\def\csname PY@tok@cpf\endcsname{\let\PY@it=\textit\def\PY@tc##1{\textcolor[rgb]{0.25,0.50,0.50}{##1}}}
\expandafter\def\csname PY@tok@c1\endcsname{\let\PY@it=\textit\def\PY@tc##1{\textcolor[rgb]{0.25,0.50,0.50}{##1}}}
\expandafter\def\csname PY@tok@cs\endcsname{\let\PY@it=\textit\def\PY@tc##1{\textcolor[rgb]{0.25,0.50,0.50}{##1}}}

\def\PYZbs{\char`\\}
\def\PYZus{\char`\_}
\def\PYZob{\char`\{}
\def\PYZcb{\char`\}}
\def\PYZca{\char`\^}
\def\PYZam{\char`\&}
\def\PYZlt{\char`\<}
\def\PYZgt{\char`\>}
\def\PYZsh{\char`\#}
\def\PYZpc{\char`\%}
\def\PYZdl{\char`\$}
\def\PYZhy{\char`\-}
\def\PYZsq{\char`\'}
\def\PYZdq{\char`\"}
\def\PYZti{\char`\~}
% for compatibility with earlier versions
\def\PYZat{@}
\def\PYZlb{[}
\def\PYZrb{]}
\makeatother


    % Exact colors from NB
    \definecolor{incolor}{rgb}{0.0, 0.0, 0.5}
    \definecolor{outcolor}{rgb}{0.545, 0.0, 0.0}



    
    % Prevent overflowing lines due to hard-to-break entities
    \sloppy 
    % Setup hyperref package
    \hypersetup{
      breaklinks=true,  % so long urls are correctly broken across lines
      colorlinks=true,
      urlcolor=urlcolor,
      linkcolor=linkcolor,
      citecolor=citecolor,
      }
    % Slightly bigger margins than the latex defaults
    
    \geometry{verbose,tmargin=1in,bmargin=1in,lmargin=1in,rmargin=1in}
    
    

    \begin{document}
    
    
    \maketitle
    
    

    
    Notas de aula: Mecânica Quântica; Autor: Jonas Maziero, UFSM

    \begin{Verbatim}[commandchars=\\\{\}]
{\color{incolor}In [{\color{incolor} }]:} \PY{k+kn}{import} \PY{n+nn}{numpy} \PY{k}{as} \PY{n+nn}{np}
        \PY{k+kn}{from} \PY{n+nn}{matplotlib} \PY{k}{import} \PY{n}{pyplot} \PY{k}{as} \PY{n}{plt}
        \PY{k+kn}{from} \PY{n+nn}{sympy} \PY{k}{import} \PY{o}{*}  \PY{c+c1}{\PYZsh{} importa tudo do sympy}
        \PY{n}{init\PYZus{}printing}\PY{p}{(}\PY{n}{use\PYZus{}unicode}\PY{o}{=}\PY{k+kc}{True}\PY{p}{)}  \PY{c+c1}{\PYZsh{} para as equações ficarem mais apresentáveis}
\end{Verbatim}


    \section{Probabilidade e
estatística}\label{probabilidade-e-estatuxedstica}

Quando não podemos prever com certeza os resultados de medidas de um
certo observável, o tratamos como uma variável aleatória, que denotamos
e.g. por X. O estado de X é descrito completamente pela distribuição de
probabilidades (DP) de seus possíveis valores. A descrição estatística
de X é obtida a partir de sua DP.

    \subsection{Variáveis discretas}\label{variuxe1veis-discretas}

\subsubsection{Probabilidade}\label{probabilidade}

Consideremos uma variável aleatória \(X\) que pode assumir os valores
\(\{X_{1},X_{2},\cdots,X_{d}\}\). A probabilidade do valor \(X_{j}\) é a
frequência relativa com que esse valor é obtido ao medirmos \(X\) um
número muito grande (idealmente \(\infty\)) de vezes:

\begin{equation}
Pr(X_{j}) = p_{j} = \frac{\text{No. de vezes que } X_{j} \text{ é obtido}}{\text{No. total de medidas de }X} = \frac{N(X_{j})}{N} = \frac{N_{j}}{N}.
\end{equation}

Claro, essas frequências são números reais entre \(0\) e \(1\), i.e.,

\begin{equation}
0 \le p_{j} \le 1.
\end{equation}

Além disso, se em cada medida obtemos um resultado válido, devemos ter
também

\begin{equation}
\sum_{j=1}^{d} N_{j} = N.
\end{equation}

Com isso vem que

\begin{equation}
\sum_{j=1}^{d} p_{j} = \sum_{j=1}^{d} \frac{N_{j}}{N} = \frac{\sum_{j=1}^{d} N_{j}}{N} = \frac{N}{N} = 1.
\end{equation}

As relações \(0 \le p_{j} \le 1\) e \(\sum_{j=1}^{d} p_{j} = 1\)
caracterizam qualquer distribuição de probabilidades válida. Podemos
fazer uma caracterização incompleta da distribuição de probabilidades de
\(X\) usando seu "centro" (dado pelo valor médio) e sua "largura" (dada
pelo desvio padrão). \#\#\# Valor médio O valor médio de uma variável
aleatória \(X\), que pode assumir os valores \(\{X_{j}\}_{j=1}^{d}\) com
probabilidades correspondentes \(\{p_{j}\}_{j=1}^{d}\), é definido como

\begin{equation}
\langle X \rangle = \sum_{j=1}^{d} X_{j}p_{j}.
\end{equation}

Ou seja, obtemos \(\langle X\rangle\) multiplicando cada valor de \(X\)
pela frequência relativa com que este é obtido e somando-se os
resultados. O valor médio não precisa necessariamente ser um dos valores
possíveis de \(X\), mas deve estar no intervalo entre seu menor e maior
valor. Vale mencionar também que se \(p_{j} = 1\) para um certo valor de
\(j\), então \(\langle X \rangle = X_{j}\). \textbf{Exercício:} Mostre
que o valor médio de uma variável aleatória está entre seus valores
mínimo e máximo, i.e., \(X_{min}\le \langle X\rangle\le X_{max}\).

\subsubsection{Valor médio de uma função de
X}\label{valor-muxe9dio-de-uma-funuxe7uxe3o-de-x}

Sem perda de generalidade, se desconsideramos possíveis degenerescências
da função \(f(X_{j})\) (valores iguais de \(f(X_{j})\) para diferentes
\(X_{j}\)), é óbvio que a frequência relativa com que obtemos \(X_{j}\)
e o valor correspondente de sua função \(f(X_{j})\) é a mesma. Por
conseguinte, como \(f(X)\) também é uma variável aleatória,

\begin{equation}
\langle f(X) \rangle = \sum_{j=1}^{d} f(X_{j})p_{j}.
\end{equation}

\subsubsection{Desvio padrão}\label{desvio-padruxe3o}

O desvio padrão é uma quantidade importante em Mecânica Quântica, e é
definido como

\begin{equation}
\sigma_{X}=\sqrt{\mathrm{Var}(X)},
\end{equation}

com a variância de \(X\) sendo

\begin{equation}
\mathrm{Var}(X)=\langle (X-\langle X\rangle)^{2}\rangle = \sum_{j=1}^{d} (X_{j}-\langle X\rangle)^{2} p_{j}.
\end{equation}

Note que \(\langle X\rangle\) é um valor escalar fixo, dada a DP.

\textbf{Exercício:} Verifique que
\(\mathrm{Var}(X)=\langle X^{2}\rangle - \langle X\rangle^{2}.\)
\textbf{Exercício:} Para entender a razão para partirmos da variância
(i.e., do quadrado da diferença \(X-\langle X\rangle\)), verifique que
\(\langle X-\langle X\rangle\rangle = 0\).

    \subsubsection{Exemplo: Estado térmico de Gibbs do oscilador harmônico
quântico
(OH)}\label{exemplo-estado-tuxe9rmico-de-gibbs-do-oscilador-harmuxf4nico-quuxe2ntico-oh}

As energias do OH são

\begin{equation}
E_{j} = \hbar\omega(j+1/2),
\end{equation}

com \(\hbar\) sendo a constante de Planck, \(\omega=2\pi f\) a
frequência angular e \(j = 0, 1, 2, \cdots, \infty\). Se o OH está em
equilíbrio com um reservatório de temperatura \(T\), a probabilidade da
energia \(E_{j}\) é

\begin{equation}
p_{j} = (1-e^{-\beta\hbar\omega})e^{-\beta\hbar\omega j},
\end{equation}

onde \(\beta = 1/k_{B}T\) com \(k_{B}\) sendo a constante de Planck.

\textbf{Exercício:} Verifique que \(\sum_{j=0}^{\infty}p_{j} = 1\).

A energia média do OH é

\begin{align}
\langle E \rangle & = \sum_{j=0}^{\infty} E_{j}p_{j} \\      
                  & = \sum_{j=0}^{\infty} \hbar\omega(j+1/2)(1-e^{-\beta\hbar\omega})e^{-\beta\hbar\omega j} \\
                  & = 2^{-1}\hbar\omega(1-e^{-\beta\hbar\omega})\sum_{j=0}^{\infty}e^{-\beta\hbar\omega j} +  \hbar\omega(1-e^{-\beta\hbar\omega})\sum_{j=0}^{\infty}je^{-\beta\hbar\omega j} \\
                  & = 2^{-1}\hbar\omega(1-e^{-\beta\hbar\omega})\sum_{j=0}^{\infty}(e^{-\beta\hbar\omega})^{j} +  \hbar\omega(1-e^{-\beta\hbar\omega})\sum_{j=0}^{\infty}j(e^{-\beta\hbar\omega})^{j}.
\end{align}

    Para continuar com esta conta precisamos de algumas séries. Partido de
\(\sum_{j=0}^{\infty}x^{j} = 1/(1-x)\) e usando
\(\partial_{x}x^{j+1} = (j+1)x^{j} = jx^{j} + x^{j}\) obteremos

\begin{align}
\sum_{j=0}^{\infty} jx^{j} & = \sum_{j=0}^{\infty} (\partial_{x}x^{j+1} - x^{j}) \\
                           & = \frac{\partial}{\partial_{x}}\left(x\sum_{j=0}^{\infty} x^{j}\right) - \sum_{j=0}^{\infty} x^{j} \\
                           & = \frac{\partial}{\partial_{x}}\left(x\frac{1}{1-x}\right) - \frac{1}{1-x} \\
                           & = \frac{1}{1-x} + \frac{x}{(1-x)^{2}} - \frac{1}{1-x} \\
                           & = \frac{x}{(1-x)^{2}}.
\end{align}

    \paragraph{Séries no sympy}\label{suxe9ries-no-sympy}

    \begin{Verbatim}[commandchars=\\\{\}]
{\color{incolor}In [{\color{incolor} }]:} \PY{n}{x}\PY{p}{,} \PY{n}{j} \PY{o}{=} \PY{n}{symbols}\PY{p}{(}\PY{l+s+s1}{\PYZsq{}}\PY{l+s+s1}{x j}\PY{l+s+s1}{\PYZsq{}}\PY{p}{)}
        \PY{n}{Sum}\PY{p}{(}\PY{n}{x}\PY{o}{*}\PY{o}{*}\PY{n}{j}\PY{p}{,}\PY{p}{(}\PY{n}{j}\PY{p}{,}\PY{l+m+mi}{0}\PY{p}{,}\PY{n}{oo}\PY{p}{)}\PY{p}{)}\PY{o}{.}\PY{n}{doit}\PY{p}{(}\PY{p}{)}
\end{Verbatim}


    Note que a condição é satisfeita, pois
\(e^{-\beta\hbar\omega} \in (0,1)\)

    \begin{Verbatim}[commandchars=\\\{\}]
{\color{incolor}In [{\color{incolor} }]:} \PY{n}{Sum}\PY{p}{(}\PY{n}{j}\PY{o}{*}\PY{n}{x}\PY{o}{*}\PY{o}{*}\PY{n}{j}\PY{p}{,}\PY{p}{(}\PY{n}{j}\PY{p}{,}\PY{l+m+mi}{0}\PY{p}{,}\PY{n}{oo}\PY{p}{)}\PY{p}{)}\PY{o}{.}\PY{n}{doit}\PY{p}{(}\PY{p}{)}
\end{Verbatim}


    Assim

\begin{align}
\langle E \rangle & = 2^{-1}\hbar\omega(1-e^{-\beta\hbar\omega})\frac{1}{1-e^{-\beta\hbar\omega}} +  \hbar\omega(1-e^{-\beta\hbar\omega})\frac{e^{-\beta\hbar\omega}}{(1-e^{-\beta\hbar\omega})^{2}} \\
                  & = \hbar\omega\left(\frac{1}{2} + \frac{e^{-\beta\hbar\omega}}{1-e^{-\beta\hbar\omega}}\right) \\
                  & = \hbar\omega\left(\frac{1}{2} + \frac{1}{e^{\beta\hbar\omega}-1}\right).
\end{align}

    \begin{Verbatim}[commandchars=\\\{\}]
{\color{incolor}In [{\color{incolor} }]:} \PY{c+c1}{\PYZsh{} não são todas as contas são facilmente feitas com a ajuda do software}
        \PY{c+c1}{\PYZsh{}b, hw = symbols(\PYZsq{}b hw\PYZsq{})}
        \PY{n}{x} \PY{o}{=} \PY{n}{symbols}\PY{p}{(}\PY{l+s+s1}{\PYZsq{}}\PY{l+s+s1}{x}\PY{l+s+s1}{\PYZsq{}}\PY{p}{,} \PY{n}{positive}\PY{o}{=}\PY{k+kc}{True}\PY{p}{)}
        \PY{k}{def} \PY{n+nf}{E}\PY{p}{(}\PY{n}{j}\PY{p}{)}\PY{p}{:}
            \PY{k}{return} \PY{n}{hw}\PY{o}{*}\PY{p}{(}\PY{l+m+mi}{1}\PY{o}{/}\PY{l+m+mi}{2}\PY{o}{+}\PY{n}{j}\PY{p}{)}
            \PY{c+c1}{\PYZsh{}return (1/2+j)}
        \PY{k}{def} \PY{n+nf}{p}\PY{p}{(}\PY{n}{j}\PY{p}{)}\PY{p}{:}
            \PY{k}{return} \PY{p}{(}\PY{l+m+mi}{1}\PY{o}{\PYZhy{}}\PY{n}{exp}\PY{p}{(}\PY{o}{\PYZhy{}}\PY{n}{b}\PY{o}{*}\PY{n}{hw}\PY{p}{)}\PY{p}{)}\PY{o}{*}\PY{n}{exp}\PY{p}{(}\PY{o}{\PYZhy{}}\PY{n}{b}\PY{o}{*}\PY{n}{hw}\PY{o}{*}\PY{n}{j}\PY{p}{)}
            \PY{c+c1}{\PYZsh{}return (1\PYZhy{}exp(\PYZhy{}x))*exp(\PYZhy{}x*j)}
        \PY{c+c1}{\PYZsh{}Sum(E(j)*p(j),(j,0,oo)).doit()}
        \PY{n}{Sum}\PY{p}{(}\PY{p}{(}\PY{l+m+mi}{1}\PY{o}{/}\PY{l+m+mi}{2}\PY{o}{+}\PY{n}{j}\PY{p}{)}\PY{o}{*}\PY{p}{(}\PY{l+m+mi}{1}\PY{o}{\PYZhy{}}\PY{n}{exp}\PY{p}{(}\PY{o}{\PYZhy{}}\PY{n}{x}\PY{p}{)}\PY{p}{)}\PY{o}{*}\PY{n}{exp}\PY{p}{(}\PY{o}{\PYZhy{}}\PY{n}{x}\PY{o}{*}\PY{n}{j}\PY{p}{)}\PY{p}{,}\PY{p}{(}\PY{n}{j}\PY{p}{,}\PY{l+m+mi}{0}\PY{p}{,}\PY{n}{oo}\PY{p}{)}\PY{p}{)}\PY{o}{.}\PY{n}{doit}\PY{p}{(}\PY{p}{)}
        \PY{c+c1}{\PYZsh{}simplify((1\PYZhy{}exp(\PYZhy{}x))*((1/2)*(Sum(exp(\PYZhy{}x*j),(j,0,oo)).doit()) + (Sum(j*exp(\PYZhy{}x*j),(j,0,oo)).doit())).doit())}
\end{Verbatim}


    \begin{Verbatim}[commandchars=\\\{\}]
{\color{incolor}In [{\color{incolor} }]:} \PY{c+c1}{\PYZsh{} Derivadas}
        \PY{n}{simplify}\PY{p}{(}\PY{n}{diff}\PY{p}{(}\PY{n}{diff}\PY{p}{(}\PY{n}{x}\PY{o}{*}\PY{o}{*}\PY{p}{(}\PY{n}{j}\PY{o}{+}\PY{l+m+mi}{2}\PY{p}{)}\PY{p}{,} \PY{n}{x}\PY{p}{)}\PY{p}{,}\PY{n}{x}\PY{p}{)}\PY{p}{)}
\end{Verbatim}


    Agora, usando
\(\partial_{xx}x^{j+2} = \partial_{x}(j+2)x^{j+1} = (j+1)(j+2)x^{j} = (j^{2} + 3j + 2)x^{j}\)
obteremos

\begin{align}
\sum_{j=0}^{\infty}j^{2}x^{j} & = \partial_{xx}\left(x^{2}\sum_{j=0}^{\infty}x^{j}\right) - 3\sum_{j=0}^{\infty}jx^{j} - 2\sum_{j=0}^{\infty}x^{j} \\
                              & = \partial_{xx}\left(x^{2}\frac{1}{1-x}\right) - 3\frac{x}{(1-x)^{2}} - 2\frac{1}{1-x} \\
                              & = \partial_{x}\left(2x\frac{1}{1-x} + x^{2}\frac{1}{(1-x)^{2}}\right) - 3\frac{x}{(1-x)^{2}} - 2\frac{1}{1-x} \\
                              & = 2\frac{1}{1-x} + 2x\frac{1}{(1-x)^{2}} + 2x\frac{1}{(1-x)^{2}} + 2x^{2}\frac{1}{(1-x)^{3}} - 3\frac{x}{(1-x)^{2}} - 2\frac{1}{1-x} \\
                              & = 2x^{2}\frac{1}{(1-x)^{3}} + x\frac{1}{(1-x)^{2}} \\
                              & = x\frac{1}{(1-x)^{2}}\left(\frac{2x}{1-x} + 1\right) \\
                              & = x\frac{1}{(1-x)^{2}}\frac{1+x}{1-x} \\
                              & = \frac{x(1+x)}{(1-x)^{3}}.
\end{align}

    \begin{Verbatim}[commandchars=\\\{\}]
{\color{incolor}In [{\color{incolor} }]:} \PY{n}{Sum}\PY{p}{(}\PY{p}{(}\PY{n}{j}\PY{o}{*}\PY{o}{*}\PY{l+m+mi}{2}\PY{p}{)}\PY{o}{*}\PY{n}{x}\PY{o}{*}\PY{o}{*}\PY{n}{j}\PY{p}{,}\PY{p}{(}\PY{n}{j}\PY{p}{,}\PY{l+m+mi}{0}\PY{p}{,}\PY{n}{oo}\PY{p}{)}\PY{p}{)}\PY{o}{.}\PY{n}{doit}\PY{p}{(}\PY{p}{)}
\end{Verbatim}


    Com isso vem que

\begin{align}
\langle E^{2}\rangle & = \sum_{j=0}^{\infty} E_{j}^{2}p_{j} \\
                     & = \sum_{j=0}^{\infty} \hbar^{2}\omega^{2}(1/2+j)^{2}(1-e^{-\beta\hbar\omega})(e^{-\beta\hbar\omega})^{j} \\
                     & = \hbar^{2}\omega^{2}(1-e^{-\beta\hbar\omega})\sum_{j=0}^{\infty} (1/4+2j+j^{2})(e^{-\beta\hbar\omega})^{j} \\
                     & = \hbar^{2}\omega^{2}(1-e^{-\beta\hbar\omega}) \left( \frac{1}{4}\sum_{j=0}^{\infty}(e^{-\beta\hbar\omega})^{j} + \sum_{j=0}^{\infty}j(e^{-\beta\hbar\omega})^{j} + \sum_{j=0}^{\infty}j^{2}(e^{-\beta\hbar\omega})^{j} \right) \\
                     & = \hbar^{2}\omega^{2}(1-e^{-\beta\hbar\omega}) \left( \frac{1}{4}\frac{1}{1-e^{-\beta\hbar\omega}} + \frac{e^{-\beta\hbar\omega}}{\left(1-e^{-\beta\hbar\omega}\right)^{2}} + \frac{e^{-\beta\hbar\omega}(1+e^{-\beta\hbar\omega})}{(1-e^{-\beta\hbar\omega})^{3}} \right) \\
                     & = \hbar^{2}\omega^{2} \left( \frac{1}{4} + \frac{e^{-\beta\hbar\omega}}{\left(1-e^{-\beta\hbar\omega}\right)} + \frac{e^{-\beta\hbar\omega}(1+e^{-\beta\hbar\omega})}{(1-e^{-\beta\hbar\omega})^{2}} \right) \\
                     & = \hbar^{2}\omega^{2} \left( \frac{1}{4} + \frac{1}{\left(e^{\beta\hbar\omega}-1\right)} + \frac{e^{\beta\hbar\omega}+1}{(e^{\beta\hbar\omega}-1)^{2}} \right).
\end{align}

Finalmente, o desvio padrão da energia é dado por

\begin{align}
\sigma_{E} & = \sqrt{\langle E^{2}\rangle - \langle E\rangle^{2}} \\
           & = \sqrt{\hbar^{2}\omega^{2} \left( \frac{1}{4} + \frac{1}{\left(e^{\beta\hbar\omega}-1\right)} + \frac{e^{\beta\hbar\omega}+1}{(e^{\beta\hbar\omega}-1)^{2}} \right) - \hbar^{2}\omega^{2}\left(\frac{1}{2} + \frac{1}{e^{\beta\hbar\omega}-1}\right)^{2}} \\
           & = \hbar\omega\sqrt{ \left( \frac{1}{4} + \frac{1}{\left(e^{\beta\hbar\omega}-1\right)} + \frac{e^{\beta\hbar\omega}+1}{(e^{\beta\hbar\omega}-1)^{2}} \right) - \left(\frac{1}{4} + \frac{1}{e^{\beta\hbar\omega}-1} + \frac{1}{(e^{\beta\hbar\omega}-1)^{2}}\right)} \\
           & = \hbar\omega\sqrt{\frac{e^{\beta\hbar\omega}}{(e^{\beta\hbar\omega}-1)^{2}}} \\
           & = \hbar\omega\frac{e^{\beta\hbar\omega/2}}{e^{\beta\hbar\omega}-1}.
\end{align}

\textbf{Exercício:} Mostre que \(\sigma_{E} = \hbar\omega\sigma_{j}\) e
que \(\sigma_{j} = \sqrt{\langle j\rangle(\langle j\rangle +1)}\).

    \begin{Verbatim}[commandchars=\\\{\}]
{\color{incolor}In [{\color{incolor} }]:} \PY{n+nb}{print}\PY{p}{(}\PY{l+s+s1}{\PYZsq{}}\PY{l+s+s1}{Distribuição de probabilidades do OH quântico em equilíbrio térmico}\PY{l+s+s1}{\PYZsq{}}\PY{p}{)}
        \PY{k+kn}{import} \PY{n+nn}{numpy} \PY{k}{as} \PY{n+nn}{np}
        \PY{k+kn}{from} \PY{n+nn}{matplotlib} \PY{k}{import} \PY{n}{pyplot} \PY{k}{as} \PY{n}{plt}
        \PY{k+kn}{from} \PY{n+nn}{matplotlib} \PY{k}{import} \PY{n}{animation}
        \PY{o}{\PYZpc{}}\PY{k}{matplotlib} nbagg
        \PY{c+c1}{\PYZsh{}\PYZpc{}matplotlib inline}
        \PY{n}{fig} \PY{o}{=} \PY{n}{plt}\PY{o}{.}\PY{n}{figure}\PY{p}{(}\PY{p}{)}
        \PY{n}{ax} \PY{o}{=} \PY{n}{plt}\PY{o}{.}\PY{n}{axes}\PY{p}{(}\PY{n}{xlim}\PY{o}{=}\PY{p}{(}\PY{l+m+mi}{0}\PY{p}{,} \PY{l+m+mi}{20}\PY{p}{)}\PY{p}{,} \PY{n}{ylim}\PY{o}{=}\PY{p}{(}\PY{l+m+mi}{0}\PY{p}{,} \PY{l+m+mi}{1}\PY{p}{)}\PY{p}{)}
        \PY{n}{line}\PY{p}{,} \PY{o}{=} \PY{n}{ax}\PY{o}{.}\PY{n}{plot}\PY{p}{(}\PY{p}{[}\PY{p}{]}\PY{p}{,} \PY{p}{[}\PY{p}{]}\PY{p}{,} \PY{n}{lw}\PY{o}{=}\PY{l+m+mi}{2}\PY{p}{)}
        
        \PY{k}{def} \PY{n+nf}{init}\PY{p}{(}\PY{p}{)}\PY{p}{:}
            \PY{n}{line}\PY{o}{.}\PY{n}{set\PYZus{}data}\PY{p}{(}\PY{p}{[}\PY{p}{]}\PY{p}{,} \PY{p}{[}\PY{p}{]}\PY{p}{)}
            \PY{k}{return} \PY{n}{line}\PY{p}{,}
        
        \PY{k}{def} \PY{n+nf}{animate}\PY{p}{(}\PY{n}{i}\PY{p}{)}\PY{p}{:}
            \PY{n}{x} \PY{o}{=} \PY{n}{np}\PY{o}{.}\PY{n}{linspace}\PY{p}{(}\PY{l+m+mi}{0}\PY{p}{,} \PY{l+m+mi}{20}\PY{p}{,} \PY{l+m+mi}{20}\PY{p}{)}
            \PY{n}{y} \PY{o}{=} \PY{p}{(}\PY{l+m+mi}{1}\PY{o}{\PYZhy{}}\PY{n}{np}\PY{o}{.}\PY{n}{exp}\PY{p}{(}\PY{o}{\PYZhy{}}\PY{l+m+mf}{0.01}\PY{o}{*}\PY{n}{i}\PY{p}{)}\PY{p}{)}\PY{o}{*}\PY{p}{(}\PY{n}{np}\PY{o}{.}\PY{n}{exp}\PY{p}{(}\PY{o}{\PYZhy{}}\PY{l+m+mf}{0.01}\PY{o}{*}\PY{n}{i}\PY{p}{)}\PY{p}{)}\PY{o}{*}\PY{o}{*}\PY{n}{x}
            \PY{n}{line}\PY{o}{.}\PY{n}{set\PYZus{}data}\PY{p}{(}\PY{n}{x}\PY{p}{,} \PY{n}{y}\PY{p}{)}
            \PY{k}{return} \PY{n}{line}\PY{p}{,}
        
        \PY{n}{ani} \PY{o}{=} \PY{n}{animation}\PY{o}{.}\PY{n}{FuncAnimation}\PY{p}{(}\PY{n}{fig}\PY{p}{,} \PY{n}{animate}\PY{p}{,} \PY{n}{init\PYZus{}func}\PY{o}{=}\PY{n}{init}\PY{p}{,} \PY{n}{frames}\PY{o}{=}\PY{l+m+mi}{300}\PY{p}{,}
                                      \PY{n}{interval}\PY{o}{=}\PY{l+m+mi}{50}\PY{p}{,} \PY{n}{blit}\PY{o}{=}\PY{k+kc}{True}\PY{p}{)}
        \PY{n}{plt}\PY{o}{.}\PY{n}{show}\PY{p}{(}\PY{p}{)}
\end{Verbatim}


    \begin{Verbatim}[commandchars=\\\{\}]
{\color{incolor}In [{\color{incolor} }]:} \PY{o}{\PYZpc{}}\PY{k}{matplotlib} inline
        \PY{k+kn}{import} \PY{n+nn}{numpy} \PY{k}{as} \PY{n+nn}{np}
        \PY{k+kn}{from} \PY{n+nn}{matplotlib} \PY{k}{import} \PY{n}{pyplot} \PY{k}{as} \PY{n}{plt}
        \PY{n}{bho} \PY{o}{=} \PY{l+m+mi}{1}
        \PY{n}{N} \PY{o}{=} \PY{l+m+mi}{20}
        \PY{n}{x} \PY{o}{=} \PY{n}{np}\PY{o}{.}\PY{n}{zeros}\PY{p}{(}\PY{n}{N}\PY{p}{)}
        \PY{n}{y} \PY{o}{=} \PY{n}{np}\PY{o}{.}\PY{n}{zeros}\PY{p}{(}\PY{n}{N}\PY{p}{)}
        \PY{k}{for} \PY{n}{j} \PY{o+ow}{in} \PY{n+nb}{range}\PY{p}{(}\PY{l+m+mi}{0}\PY{p}{,}\PY{n}{N}\PY{p}{)}\PY{p}{:}
            \PY{n}{x}\PY{p}{[}\PY{n}{j}\PY{p}{]} \PY{o}{=} \PY{n}{j}
            \PY{n}{y}\PY{p}{[}\PY{n}{j}\PY{p}{]} \PY{o}{=} \PY{p}{(}\PY{l+m+mi}{1}\PY{o}{\PYZhy{}}\PY{n}{np}\PY{o}{.}\PY{n}{exp}\PY{p}{(}\PY{o}{\PYZhy{}}\PY{n}{bho}\PY{p}{)}\PY{p}{)}\PY{o}{*}\PY{p}{(}\PY{n}{np}\PY{o}{.}\PY{n}{exp}\PY{p}{(}\PY{o}{\PYZhy{}}\PY{n}{bho}\PY{p}{)}\PY{p}{)}\PY{o}{*}\PY{o}{*}\PY{n}{j}
        \PY{n}{plt}\PY{o}{.}\PY{n}{bar}\PY{p}{(}\PY{n}{x}\PY{p}{,} \PY{n}{y}\PY{p}{,} \PY{n}{label} \PY{o}{=} \PY{l+s+s1}{\PYZsq{}}\PY{l+s+s1}{\PYZsq{}}\PY{p}{,} \PY{n}{color} \PY{o}{=} \PY{l+s+s1}{\PYZsq{}}\PY{l+s+s1}{red}\PY{l+s+s1}{\PYZsq{}}\PY{p}{)}
        \PY{n}{plt}\PY{o}{.}\PY{n}{xlabel}\PY{p}{(}\PY{l+s+s1}{\PYZsq{}}\PY{l+s+s1}{j}\PY{l+s+s1}{\PYZsq{}}\PY{p}{)}
        \PY{n}{plt}\PY{o}{.}\PY{n}{ylabel}\PY{p}{(}\PY{l+s+s1}{\PYZsq{}}\PY{l+s+s1}{pj}\PY{l+s+s1}{\PYZsq{}}\PY{p}{)}
        \PY{c+c1}{\PYZsh{}plt.legend()}
        \PY{n}{axes} \PY{o}{=} \PY{n}{plt}\PY{o}{.}\PY{n}{gca}\PY{p}{(}\PY{p}{)}
        \PY{n}{axes}\PY{o}{.}\PY{n}{set\PYZus{}xlim}\PY{p}{(}\PY{p}{[}\PY{l+m+mi}{0}\PY{o}{\PYZhy{}}\PY{l+m+mf}{0.5}\PY{p}{,}\PY{n}{N}\PY{p}{]}\PY{p}{)}
        \PY{n}{axes}\PY{o}{.}\PY{n}{set\PYZus{}ylim}\PY{p}{(}\PY{p}{[}\PY{o}{\PYZhy{}}\PY{l+m+mf}{0.1}\PY{p}{,}\PY{l+m+mi}{1}\PY{p}{]}\PY{p}{)}
        \PY{n}{plt}\PY{o}{.}\PY{n}{show}\PY{p}{(}\PY{p}{)}
\end{Verbatim}


    \subsection{Variáveis contínuas}\label{variuxe1veis-contuxednuas}

Neste caso, lidaremos com observáveis que possuem um espectro contínuo,
i.e., seu conjunto de valores possível não é enumerável como no exemplo
acima.

\subsubsection{Densidade de
probabilidade}\label{densidade-de-probabilidade}

Como não faz sentido falar na probabilidade de um valor específico de
\(X\), usamos a densidade de probabilidade \(\rho(x)\), que é a
probabilidade por intervalo infinitesimal \(dx\) de valores de \(X\).
Teremos assim que

\begin{equation}
\rho(x)dx = \text{probabilidade de obter o valor de } X \text{ entre } x \text{ e } x+dx.
\end{equation}

Se conhecemos a função \(\rho(x)\), para calcular a frequência relativa
com que valores de \(X\) são obtidos dentro de um certo intervalo finito
\([a,b]\) simplesmente somamos as probabilidades nos sub-intervalos
infinitesimais aí:

\begin{equation}
Pr(x\in[a,b]) = \int_{a}^{b} \rho(x)dx.
\end{equation}

Analogamente ao caso discreto, teremos a seguinte condição de
normalização para a distribuição de probabilidades

\begin{equation}
\int_{-\infty}^{+\infty} \rho(x)dx = 1.
\end{equation}

\subsubsection{Valor médio e desvio
padrão}\label{valor-muxe9dio-e-desvio-padruxe3o}

O cálculo dessas quantidades para variáveis contínuas segue a mesma idea
que para variáveis discretas, ou seja,

\begin{align}
\langle X \rangle & = \int_{-\infty}^{+\infty} x\rho(x)dx, \\
\langle f(X) \rangle & = \int_{-\infty}^{+\infty} f(x)\rho(x)dx, \\
\sigma_{X}^{2} & = \langle (X - \langle X \rangle)^{2} \rangle = \int_{-\infty}^{+\infty} (x - \langle X \rangle)^{2}\rho(x)dx = \langle X^{2} \rangle - \langle X \rangle^{2}. \\
\end{align}

    \subsubsection{Exemplo: Distribuição Gaussiana de
velocidades}\label{exemplo-distribuiuxe7uxe3o-gaussiana-de-velocidades}

Esta é uma das DP mais importantes para a ciência e tem a forma

\begin{equation}
\rho_{G}(x) = \frac{1}{\sqrt{2\pi \sigma^{2}}}e^{-(x-\mu)^{2}/2\sigma^{2}}.
\end{equation}

    \paragraph{\texorpdfstring{Integral Gaussiana
\(G = \int_{-\infty}^{\infty}e^{-ax^{2}}dx = \sqrt{\frac{\pi}{a}}\)}{Integral Gaussiana G = \textbackslash{}int\_\{-\textbackslash{}infty\}\^{}\{\textbackslash{}infty\}e\^{}\{-ax\^{}\{2\}\}dx = \textbackslash{}sqrt\{\textbackslash{}frac\{\textbackslash{}pi\}\{a\}\}}}\label{integral-gaussiana-g-int_-inftyinftye-ax2dx-sqrtfracpia}

Fazemos essa integral notando que

\begin{align}
G^{2} & = \left(\int_{-\infty}^{\infty}e^{-ax^{2}}dx\right)\left(\int_{-\infty}^{\infty}e^{-ay^{2}}dy\right) \\
      & = \int_{-\infty}^{\infty}\int_{-\infty}^{\infty} e^{-a(x^{2}+y^{2})} dx dy.
\end{align}

Agora mudamos de coordenadas retangulares para polares:
\(x = r\cos\theta\), \(y = r\sin\theta\), \(da = dxdy = rdrd\theta\) com
\(x\in(-\infty,\infty)\), \(y\in(-\infty,\infty)\),
\(\theta\in[0,\infty)\), \(\theta\in[0,2\pi)\). Assim

\begin{align}
G^{2} & = \int_{0}^{\infty}\int_{0}^{2\pi} e^{-ar^2} rdrd\theta \\
      & = 2\pi\int_{0}^{-\infty} e^{u} \frac{-du}{2a} \\
      & = \frac{\pi}{a}.
\end{align}

    \begin{Verbatim}[commandchars=\\\{\}]
{\color{incolor}In [{\color{incolor} }]:} \PY{n}{x}\PY{p}{,} \PY{n}{a} \PY{o}{=} \PY{n}{symbols}\PY{p}{(}\PY{l+s+s1}{\PYZsq{}}\PY{l+s+s1}{x a}\PY{l+s+s1}{\PYZsq{}}\PY{p}{)}
        \PY{n}{integrate}\PY{p}{(}\PY{n}{exp}\PY{p}{(}\PY{o}{\PYZhy{}}\PY{n}{a}\PY{o}{*}\PY{n}{x}\PY{o}{*}\PY{o}{*}\PY{l+m+mi}{2}\PY{p}{)}\PY{p}{,}\PY{p}{(}\PY{n}{x}\PY{p}{,}\PY{o}{\PYZhy{}}\PY{n}{oo}\PY{p}{,}\PY{n}{oo}\PY{p}{)}\PY{p}{)}
\end{Verbatim}


    \textbf{Exercício:} Verifique que
\(\int_{-\infty}^{\infty}\rho_{G}(x)dx = 1\).

    Com isso já podemos calcular o valor médio de \(X\) com DP Gaussiana:

\begin{align}
\langle X \rangle & = \int_{-\infty}^{\infty} x\rho_{G}(x)dx \\
                  & = \int_{-\infty}^{\infty} x\frac{1}{\sqrt{2\pi \sigma^{2}}}e^{-(x-\mu)^{2}/2\sigma^{2}}dx.
\end{align}

Fazendo a mudança de variável
\(y = (x-\mu)/\sqrt{2}\sigma \therefore dy=dx/\sqrt{2}\sigma\) obtemos
(para \(\mu\) finito)

\begin{align}
\langle X \rangle & = \frac{1}{\sqrt{2\pi \sigma^{2}}}\int_{-\infty}^{\infty} (\sqrt{2}\sigma y+\mu)e^{-y^{2}} \sqrt{2} \sigma dy \\
                  & = \frac{1}{\sqrt{2\pi \sigma^{2}}}\left (2\sigma^{2}\int_{-\infty}^{\infty}y e^{-y^{2}}dy + \sqrt{2}\sigma\mu \int_{-\infty}^{\infty} e^{-y^{2}}dy \right) \\
                  & = \frac{1}{\sqrt{2\pi \sigma^{2}}}\left (2\sigma^{2}0 + \sqrt{2}\sigma\mu\sqrt{\pi/1} \right) \\
                  & = \mu.
\end{align}

Acima a primeira integral é nula pois envolve uma função ímpar.

    \paragraph{Integrais relacionadas à Gaussiana obtidas usando o truque de
Feynman}\label{integrais-relacionadas-uxe0-gaussiana-obtidas-usando-o-truque-de-feynman}

Primeiro notemos que \(\partial_{a}e^{-a x^{2}} = -x^{2}e^{-a x^{2}}\).

    \begin{Verbatim}[commandchars=\\\{\}]
{\color{incolor}In [{\color{incolor} }]:} \PY{n}{diff}\PY{p}{(}\PY{n}{exp}\PY{p}{(}\PY{o}{\PYZhy{}}\PY{n}{a}\PY{o}{*}\PY{n}{x}\PY{o}{*}\PY{o}{*}\PY{l+m+mi}{2}\PY{p}{)}\PY{p}{,}\PY{n}{a}\PY{p}{)}
\end{Verbatim}


    Isso pode ser usado para a integral

\begin{align}
\int_{-\infty}^{\infty} x^{2}e^{-ax^{2}}dx & = \int_{-\infty}^{\infty} (-\partial_{a}e^{-a x^{2}})dx \\
                                           & = -\partial_{a}\int_{-\infty}^{\infty} e^{-a x^{2}}dx \\
                                           & = -\partial_{a}\sqrt{\pi/a} \\
                                           & = \sqrt{\pi/4 a^{3}}.
\end{align}

    \begin{Verbatim}[commandchars=\\\{\}]
{\color{incolor}In [{\color{incolor} }]:} \PY{n}{x}\PY{p}{,} \PY{n}{a} \PY{o}{=} \PY{n}{symbols}\PY{p}{(}\PY{l+s+s1}{\PYZsq{}}\PY{l+s+s1}{x a}\PY{l+s+s1}{\PYZsq{}}\PY{p}{)}
        \PY{n}{integrate}\PY{p}{(}\PY{p}{(}\PY{n}{x}\PY{o}{*}\PY{o}{*}\PY{l+m+mi}{2}\PY{p}{)}\PY{o}{*}\PY{n}{exp}\PY{p}{(}\PY{o}{\PYZhy{}}\PY{n}{a}\PY{o}{*}\PY{n}{x}\PY{o}{*}\PY{o}{*}\PY{l+m+mi}{2}\PY{p}{)}\PY{p}{,}\PY{p}{(}\PY{n}{x}\PY{p}{,}\PY{o}{\PYZhy{}}\PY{n}{oo}\PY{p}{,}\PY{n}{oo}\PY{p}{)}\PY{p}{)}
\end{Verbatim}


    Teremos assim que

\begin{align}
\langle X^{2} \rangle & = \int_{-\infty}^{\infty} x^{2} \rho_{G}(x)dx \\
                  & = \frac{1}{\sqrt{2\pi \sigma^{2}}} \int_{-\infty}^{\infty} x^{2} e^{-(x-\mu)^{2}/2\sigma^{2}}dx \\
                  & = \frac{1}{\sqrt{2\pi \sigma^{2}}} \int_{-\infty}^{\infty} \left(\sqrt{2}\sigma y+\mu\right)^{2} e^{-y^{2}}\sqrt{2}\sigma dy \\
                  & = \frac{1}{\sqrt{\pi}} \int_{-\infty}^{\infty} \left(2\sigma^{2}y^{2}+2\sqrt{2}\sigma\mu y+\mu^{2}\right) e^{-y^{2}}dy \\
                  & = \frac{1}{\sqrt{\pi}} \left(2\sigma^{2}\int_{-\infty}^{\infty} y^{2} e^{-y^{2}}dy + 2\sqrt{2}\sigma\mu\int_{-\infty}^{\infty}  y e^{-y^{2}}dy + \mu^{2}\int_{-\infty}^{\infty}  e^{-y^{2}}dy \right) \\
                  & = \frac{1}{\sqrt{\pi}} \left(2\sigma^{2}\sqrt{\pi/4} + 2\sqrt{2}\sigma\mu0 + \mu^{2}\sqrt{\pi} \right) \\
                  & = \sigma^{2} + \mu^{2}.
\end{align}

\paragraph{'Desvio padrão'}\label{desvio-padruxe3o}

\begin{align}
\sigma_{X}^{2} & = \langle X^{2} \rangle - \langle X \rangle^{2} \\
           & = \sigma^{2} + \mu^{2} - \mu^{2} \\
           & = \sigma^{2}.
\end{align}

    \begin{Verbatim}[commandchars=\\\{\}]
{\color{incolor}In [{\color{incolor} }]:} \PY{n+nb}{print}\PY{p}{(}\PY{l+s+s1}{\PYZsq{}}\PY{l+s+s1}{Distribuição de probabilidades Gaussiana}\PY{l+s+s1}{\PYZsq{}}\PY{p}{)}
        \PY{o}{\PYZpc{}}\PY{k}{matplotlib} inline
        \PY{k+kn}{import} \PY{n+nn}{numpy} \PY{k}{as} \PY{n+nn}{np}
        \PY{k+kn}{from} \PY{n+nn}{matplotlib} \PY{k}{import} \PY{n}{pyplot} \PY{k}{as} \PY{n}{plt}
        \PY{k+kn}{from} \PY{n+nn}{math} \PY{k}{import} \PY{n}{sqrt}\PY{p}{,} \PY{n}{exp}\PY{p}{,} \PY{n}{pi}
        \PY{n}{mu1} \PY{o}{=} \PY{l+m+mi}{0}
        \PY{n}{sd1} \PY{o}{=} \PY{l+m+mf}{0.3}
        \PY{n}{mu2} \PY{o}{=} \PY{l+m+mi}{0}
        \PY{n}{sd2} \PY{o}{=} \PY{l+m+mi}{1}
        \PY{n}{mu3} \PY{o}{=} \PY{l+m+mi}{0}
        \PY{n}{sd3} \PY{o}{=} \PY{l+m+mi}{2}
        \PY{n}{mu4} \PY{o}{=} \PY{l+m+mi}{3}
        \PY{n}{sd4} \PY{o}{=} \PY{l+m+mi}{1}
        \PY{n}{mu5} \PY{o}{=} \PY{o}{\PYZhy{}}\PY{l+m+mi}{3}
        \PY{n}{sd5} \PY{o}{=} \PY{l+m+mi}{1}
        \PY{n}{N} \PY{o}{=} \PY{l+m+mi}{60}
        \PY{n}{x} \PY{o}{=} \PY{n}{np}\PY{o}{.}\PY{n}{zeros}\PY{p}{(}\PY{l+m+mi}{2}\PY{o}{*}\PY{n}{N}\PY{p}{)}
        \PY{n}{y1} \PY{o}{=} \PY{n}{np}\PY{o}{.}\PY{n}{zeros}\PY{p}{(}\PY{l+m+mi}{2}\PY{o}{*}\PY{n}{N}\PY{p}{)}
        \PY{n}{y2} \PY{o}{=} \PY{n}{np}\PY{o}{.}\PY{n}{zeros}\PY{p}{(}\PY{l+m+mi}{2}\PY{o}{*}\PY{n}{N}\PY{p}{)}
        \PY{n}{y3} \PY{o}{=} \PY{n}{np}\PY{o}{.}\PY{n}{zeros}\PY{p}{(}\PY{l+m+mi}{2}\PY{o}{*}\PY{n}{N}\PY{p}{)}
        \PY{n}{y4} \PY{o}{=} \PY{n}{np}\PY{o}{.}\PY{n}{zeros}\PY{p}{(}\PY{l+m+mi}{2}\PY{o}{*}\PY{n}{N}\PY{p}{)}
        \PY{n}{y5} \PY{o}{=} \PY{n}{np}\PY{o}{.}\PY{n}{zeros}\PY{p}{(}\PY{l+m+mi}{2}\PY{o}{*}\PY{n}{N}\PY{p}{)}
        \PY{k}{for} \PY{n}{j} \PY{o+ow}{in} \PY{n+nb}{range}\PY{p}{(}\PY{l+m+mi}{0}\PY{p}{,}\PY{l+m+mi}{2}\PY{o}{*}\PY{n}{N}\PY{p}{)}\PY{p}{:}
            \PY{n}{x}\PY{p}{[}\PY{n}{j}\PY{p}{]} \PY{o}{=} \PY{p}{(}\PY{n}{j}\PY{o}{\PYZhy{}}\PY{n}{N}\PY{p}{)}\PY{o}{/}\PY{l+m+mi}{6}
            \PY{n}{y1}\PY{p}{[}\PY{n}{j}\PY{p}{]} \PY{o}{=} \PY{p}{(}\PY{l+m+mi}{1}\PY{o}{/}\PY{n}{sqrt}\PY{p}{(}\PY{l+m+mi}{2}\PY{o}{*}\PY{n}{pi}\PY{o}{*}\PY{n}{sd1}\PY{o}{*}\PY{o}{*}\PY{l+m+mi}{2}\PY{p}{)}\PY{p}{)}\PY{o}{*}\PY{n}{exp}\PY{p}{(}\PY{o}{\PYZhy{}}\PY{p}{(}\PY{p}{(}\PY{n}{x}\PY{p}{[}\PY{n}{j}\PY{p}{]}\PY{o}{\PYZhy{}}\PY{n}{mu1}\PY{p}{)}\PY{o}{*}\PY{o}{*}\PY{l+m+mi}{2}\PY{p}{)}\PY{o}{/}\PY{p}{(}\PY{l+m+mi}{2}\PY{o}{*}\PY{n}{sd1}\PY{o}{*}\PY{o}{*}\PY{l+m+mi}{2}\PY{p}{)}\PY{p}{)}
            \PY{n}{y2}\PY{p}{[}\PY{n}{j}\PY{p}{]} \PY{o}{=} \PY{p}{(}\PY{l+m+mi}{1}\PY{o}{/}\PY{n}{sqrt}\PY{p}{(}\PY{l+m+mi}{2}\PY{o}{*}\PY{n}{pi}\PY{o}{*}\PY{n}{sd2}\PY{o}{*}\PY{o}{*}\PY{l+m+mi}{2}\PY{p}{)}\PY{p}{)}\PY{o}{*}\PY{n}{exp}\PY{p}{(}\PY{o}{\PYZhy{}}\PY{p}{(}\PY{p}{(}\PY{n}{x}\PY{p}{[}\PY{n}{j}\PY{p}{]}\PY{o}{\PYZhy{}}\PY{n}{mu2}\PY{p}{)}\PY{o}{*}\PY{o}{*}\PY{l+m+mi}{2}\PY{p}{)}\PY{o}{/}\PY{p}{(}\PY{l+m+mi}{2}\PY{o}{*}\PY{n}{sd2}\PY{o}{*}\PY{o}{*}\PY{l+m+mi}{2}\PY{p}{)}\PY{p}{)}
            \PY{n}{y3}\PY{p}{[}\PY{n}{j}\PY{p}{]} \PY{o}{=} \PY{p}{(}\PY{l+m+mi}{1}\PY{o}{/}\PY{n}{sqrt}\PY{p}{(}\PY{l+m+mi}{2}\PY{o}{*}\PY{n}{pi}\PY{o}{*}\PY{n}{sd3}\PY{o}{*}\PY{o}{*}\PY{l+m+mi}{2}\PY{p}{)}\PY{p}{)}\PY{o}{*}\PY{n}{exp}\PY{p}{(}\PY{o}{\PYZhy{}}\PY{p}{(}\PY{p}{(}\PY{n}{x}\PY{p}{[}\PY{n}{j}\PY{p}{]}\PY{o}{\PYZhy{}}\PY{n}{mu3}\PY{p}{)}\PY{o}{*}\PY{o}{*}\PY{l+m+mi}{2}\PY{p}{)}\PY{o}{/}\PY{p}{(}\PY{l+m+mi}{2}\PY{o}{*}\PY{n}{sd3}\PY{o}{*}\PY{o}{*}\PY{l+m+mi}{2}\PY{p}{)}\PY{p}{)}
            \PY{n}{y4}\PY{p}{[}\PY{n}{j}\PY{p}{]} \PY{o}{=} \PY{p}{(}\PY{l+m+mi}{1}\PY{o}{/}\PY{n}{sqrt}\PY{p}{(}\PY{l+m+mi}{2}\PY{o}{*}\PY{n}{pi}\PY{o}{*}\PY{n}{sd4}\PY{o}{*}\PY{o}{*}\PY{l+m+mi}{2}\PY{p}{)}\PY{p}{)}\PY{o}{*}\PY{n}{exp}\PY{p}{(}\PY{o}{\PYZhy{}}\PY{p}{(}\PY{p}{(}\PY{n}{x}\PY{p}{[}\PY{n}{j}\PY{p}{]}\PY{o}{\PYZhy{}}\PY{n}{mu4}\PY{p}{)}\PY{o}{*}\PY{o}{*}\PY{l+m+mi}{2}\PY{p}{)}\PY{o}{/}\PY{p}{(}\PY{l+m+mi}{2}\PY{o}{*}\PY{n}{sd4}\PY{o}{*}\PY{o}{*}\PY{l+m+mi}{2}\PY{p}{)}\PY{p}{)}
            \PY{n}{y5}\PY{p}{[}\PY{n}{j}\PY{p}{]} \PY{o}{=} \PY{p}{(}\PY{l+m+mi}{1}\PY{o}{/}\PY{n}{sqrt}\PY{p}{(}\PY{l+m+mi}{2}\PY{o}{*}\PY{n}{pi}\PY{o}{*}\PY{n}{sd5}\PY{o}{*}\PY{o}{*}\PY{l+m+mi}{2}\PY{p}{)}\PY{p}{)}\PY{o}{*}\PY{n}{exp}\PY{p}{(}\PY{o}{\PYZhy{}}\PY{p}{(}\PY{p}{(}\PY{n}{x}\PY{p}{[}\PY{n}{j}\PY{p}{]}\PY{o}{\PYZhy{}}\PY{n}{mu5}\PY{p}{)}\PY{o}{*}\PY{o}{*}\PY{l+m+mi}{2}\PY{p}{)}\PY{o}{/}\PY{p}{(}\PY{l+m+mi}{2}\PY{o}{*}\PY{n}{sd5}\PY{o}{*}\PY{o}{*}\PY{l+m+mi}{2}\PY{p}{)}\PY{p}{)}
        \PY{n}{plt}\PY{o}{.}\PY{n}{plot}\PY{p}{(}\PY{n}{x}\PY{p}{,} \PY{n}{y1}\PY{p}{,} \PY{n}{label} \PY{o}{=} \PY{l+s+s1}{\PYZsq{}}\PY{l+s+s1}{\PYZsq{}}\PY{p}{,} \PY{n}{color} \PY{o}{=} \PY{l+s+s1}{\PYZsq{}}\PY{l+s+s1}{blue}\PY{l+s+s1}{\PYZsq{}}\PY{p}{)}
        \PY{n}{plt}\PY{o}{.}\PY{n}{plot}\PY{p}{(}\PY{n}{x}\PY{p}{,} \PY{n}{y2}\PY{p}{,} \PY{n}{label} \PY{o}{=} \PY{l+s+s1}{\PYZsq{}}\PY{l+s+s1}{\PYZsq{}}\PY{p}{,} \PY{n}{color} \PY{o}{=} \PY{l+s+s1}{\PYZsq{}}\PY{l+s+s1}{red}\PY{l+s+s1}{\PYZsq{}}\PY{p}{)}
        \PY{n}{plt}\PY{o}{.}\PY{n}{plot}\PY{p}{(}\PY{n}{x}\PY{p}{,} \PY{n}{y3}\PY{p}{,} \PY{n}{label} \PY{o}{=} \PY{l+s+s1}{\PYZsq{}}\PY{l+s+s1}{\PYZsq{}}\PY{p}{,} \PY{n}{color} \PY{o}{=} \PY{l+s+s1}{\PYZsq{}}\PY{l+s+s1}{green}\PY{l+s+s1}{\PYZsq{}}\PY{p}{)}
        \PY{n}{plt}\PY{o}{.}\PY{n}{plot}\PY{p}{(}\PY{n}{x}\PY{p}{,} \PY{n}{y4}\PY{p}{,} \PY{n}{label} \PY{o}{=} \PY{l+s+s1}{\PYZsq{}}\PY{l+s+s1}{\PYZsq{}}\PY{p}{,} \PY{n}{color} \PY{o}{=} \PY{l+s+s1}{\PYZsq{}}\PY{l+s+s1}{yellow}\PY{l+s+s1}{\PYZsq{}}\PY{p}{)}
        \PY{n}{plt}\PY{o}{.}\PY{n}{plot}\PY{p}{(}\PY{n}{x}\PY{p}{,} \PY{n}{y5}\PY{p}{,} \PY{n}{label} \PY{o}{=} \PY{l+s+s1}{\PYZsq{}}\PY{l+s+s1}{\PYZsq{}}\PY{p}{,} \PY{n}{color} \PY{o}{=} \PY{l+s+s1}{\PYZsq{}}\PY{l+s+s1}{cyan}\PY{l+s+s1}{\PYZsq{}}\PY{p}{)}
        \PY{n}{plt}\PY{o}{.}\PY{n}{xlabel}\PY{p}{(}\PY{l+s+s1}{\PYZsq{}}\PY{l+s+s1}{x}\PY{l+s+s1}{\PYZsq{}}\PY{p}{)}
        \PY{n}{plt}\PY{o}{.}\PY{n}{ylabel}\PY{p}{(}\PY{l+s+s1}{\PYZsq{}}\PY{l+s+s1}{pr(x)}\PY{l+s+s1}{\PYZsq{}}\PY{p}{)}
        \PY{n}{plt}\PY{o}{.}\PY{n}{legend}\PY{p}{(}\PY{p}{)}
        \PY{n}{axes} \PY{o}{=} \PY{n}{plt}\PY{o}{.}\PY{n}{gca}\PY{p}{(}\PY{p}{)}
        \PY{n}{axes}\PY{o}{.}\PY{n}{set\PYZus{}xlim}\PY{p}{(}\PY{p}{[}\PY{o}{\PYZhy{}}\PY{l+m+mi}{10}\PY{p}{,}\PY{l+m+mi}{10}\PY{p}{]}\PY{p}{)}
        \PY{n}{axes}\PY{o}{.}\PY{n}{set\PYZus{}ylim}\PY{p}{(}\PY{p}{[}\PY{o}{\PYZhy{}}\PY{l+m+mf}{0.1}\PY{p}{,}\PY{p}{]}\PY{p}{)}
        \PY{n}{plt}\PY{o}{.}\PY{n}{show}\PY{p}{(}\PY{p}{)}
\end{Verbatim}


    \begin{Verbatim}[commandchars=\\\{\}]
{\color{incolor}In [{\color{incolor}22}]:} \PY{o}{\PYZpc{}}\PY{k}{matplotlib} widget
\end{Verbatim}


    \begin{Verbatim}[commandchars=\\\{\}]
Warning: Cannot change to a different GUI toolkit: widget. Using nbagg instead.

    \end{Verbatim}

    \subsubsection{Teorema do limite central
(TLC)}\label{teorema-do-limite-central-tlc}

Para qualquer variável aleatória \(X\), se a medimos um número finito de
vezes, o valor médio obtido será também uma variável aleatória. O TLC
diz que a distribuição de probabilidades de \(\langle X \rangle\) é a
distribuição normal. Vamos considerar como exemplo um dado. Vamos fazer
\(M\) experimentos, cada um com \(N\) rolagens do dado. Em cada
experimento obteremos um certo valor médio do valor do dado. Vamos então
graficar a distribuição de probabilidades desses valores médios.

    \begin{Verbatim}[commandchars=\\\{\}]
{\color{incolor}In [{\color{incolor} }]:} \PY{c+c1}{\PYZsh{} M jogos, cada um com N roladas do dado}
        \PY{k+kn}{from} \PY{n+nn}{numpy} \PY{k}{import} \PY{n}{random}\PY{p}{;}  \PY{n}{random}\PY{o}{.}\PY{n}{seed}\PY{p}{(}\PY{p}{)}
        \PY{k+kn}{from} \PY{n+nn}{scipy}\PY{n+nn}{.}\PY{n+nn}{stats} \PY{k}{import} \PY{n}{norm}
        
        \PY{n}{M} \PY{o}{=} \PY{l+m+mi}{1000}\PY{p}{;}  \PY{n}{N} \PY{o}{=} \PY{l+m+mi}{2}\PY{o}{*}\PY{o}{*}\PY{l+m+mi}{13}
        \PY{n}{avg} \PY{o}{=} \PY{n}{np}\PY{o}{.}\PY{n}{zeros}\PY{p}{(}\PY{n}{M}\PY{p}{)}
        \PY{k}{for} \PY{n}{j} \PY{o+ow}{in} \PY{n+nb}{range}\PY{p}{(}\PY{l+m+mi}{0}\PY{p}{,}\PY{n}{M}\PY{p}{)}\PY{p}{:}
            \PY{n}{p} \PY{o}{=} \PY{n}{np}\PY{o}{.}\PY{n}{zeros}\PY{p}{(}\PY{l+m+mi}{6}\PY{p}{)}
            \PY{k}{for} \PY{n}{k} \PY{o+ow}{in} \PY{n+nb}{range}\PY{p}{(}\PY{l+m+mi}{0}\PY{p}{,}\PY{n}{N}\PY{p}{)}\PY{p}{:}
                \PY{n}{rn} \PY{o}{=} \PY{n}{random}\PY{o}{.}\PY{n}{random}\PY{p}{(}\PY{p}{)}
                \PY{k}{for} \PY{n}{l} \PY{o+ow}{in} \PY{n+nb}{range}\PY{p}{(}\PY{l+m+mi}{0}\PY{p}{,}\PY{l+m+mi}{6}\PY{p}{)}\PY{p}{:}
                    \PY{k}{if} \PY{n}{rn} \PY{o}{\PYZgt{}}\PY{o}{=} \PY{n}{l}\PY{o}{/}\PY{l+m+mf}{6.} \PY{o+ow}{and} \PY{n}{rn} \PY{o}{\PYZlt{}} \PY{p}{(}\PY{n}{l}\PY{o}{+}\PY{l+m+mi}{1}\PY{p}{)}\PY{o}{/}\PY{l+m+mf}{6.}\PY{p}{:}
                        \PY{n}{p}\PY{p}{[}\PY{n}{l}\PY{p}{]} \PY{o}{+}\PY{o}{=} \PY{l+m+mi}{1}
            \PY{k}{for} \PY{n}{l} \PY{o+ow}{in} \PY{n+nb}{range}\PY{p}{(}\PY{l+m+mi}{0}\PY{p}{,}\PY{l+m+mi}{6}\PY{p}{)}\PY{p}{:}
                \PY{n}{p}\PY{p}{[}\PY{n}{l}\PY{p}{]} \PY{o}{=} \PY{n}{p}\PY{p}{[}\PY{n}{l}\PY{p}{]}\PY{o}{/}\PY{n}{N}
                \PY{n}{avg}\PY{p}{[}\PY{n}{j}\PY{p}{]} \PY{o}{+}\PY{o}{=} \PY{n}{l}\PY{o}{*}\PY{n}{p}\PY{p}{[}\PY{n}{l}\PY{p}{]}
        
        \PY{n}{bins} \PY{o}{=} \PY{p}{[}\PY{l+m+mi}{0}\PY{p}{,} \PY{l+m+mf}{0.25}\PY{p}{,} \PY{l+m+mf}{0.5}\PY{p}{,} \PY{l+m+mf}{0.75}\PY{p}{,} \PY{l+m+mi}{1}\PY{p}{,} \PY{l+m+mf}{1.25}\PY{p}{,} \PY{l+m+mf}{1.5}\PY{p}{,} \PY{l+m+mf}{1.75}\PY{p}{,} \PY{l+m+mi}{2}\PY{p}{,} \PY{l+m+mf}{2.25}\PY{p}{,} \PY{l+m+mf}{2.5}\PY{p}{,} \PY{l+m+mf}{2.75}\PY{p}{,} \PY{l+m+mi}{3}\PY{p}{,} \PY{l+m+mf}{3.25}\PY{p}{,} \PY{l+m+mf}{3.5}\PY{p}{,} \PY{l+m+mf}{3.75}\PY{p}{,} \PY{l+m+mi}{4}\PY{p}{,} \PY{l+m+mf}{4.25}\PY{p}{,} \PY{l+m+mf}{4.5}\PY{p}{,} \PY{l+m+mf}{4.75}\PY{p}{,} \PY{l+m+mi}{5}\PY{p}{]}
        \PY{n}{plt}\PY{o}{.}\PY{n}{hist}\PY{p}{(}\PY{n}{avg}\PY{p}{,} \PY{n}{bins}\PY{p}{,} \PY{n}{histtype}\PY{o}{=}\PY{l+s+s1}{\PYZsq{}}\PY{l+s+s1}{bar}\PY{l+s+s1}{\PYZsq{}}\PY{p}{,} \PY{n}{rwidth}\PY{o}{=}\PY{l+m+mf}{0.8}\PY{p}{)}
        
        \PY{n}{mu}\PY{p}{,} \PY{n}{std} \PY{o}{=} \PY{n}{norm}\PY{o}{.}\PY{n}{fit}\PY{p}{(}\PY{n}{avg}\PY{p}{)}
        \PY{n}{xmin}\PY{p}{,} \PY{n}{xmax} \PY{o}{=} \PY{n}{plt}\PY{o}{.}\PY{n}{xlim}\PY{p}{(}\PY{p}{)}
        \PY{n}{x} \PY{o}{=} \PY{n}{np}\PY{o}{.}\PY{n}{linspace}\PY{p}{(}\PY{n}{xmin}\PY{p}{,} \PY{n}{xmax}\PY{p}{,} \PY{l+m+mi}{200}\PY{p}{)}
        \PY{n}{p} \PY{o}{=} \PY{p}{(}\PY{n}{M}\PY{o}{/}\PY{l+m+mi}{4}\PY{p}{)}\PY{o}{*}\PY{n}{norm}\PY{o}{.}\PY{n}{pdf}\PY{p}{(}\PY{n}{x}\PY{p}{,} \PY{n}{mu}\PY{p}{,} \PY{n}{std}\PY{p}{)}
        \PY{n}{plt}\PY{o}{.}\PY{n}{plot}\PY{p}{(}\PY{n}{x}\PY{p}{,} \PY{n}{p}\PY{p}{,} \PY{l+s+s1}{\PYZsq{}}\PY{l+s+s1}{k}\PY{l+s+s1}{\PYZsq{}}\PY{p}{,} \PY{n}{linewidth}\PY{o}{=}\PY{l+m+mi}{2}\PY{p}{)}
        \PY{n}{title} \PY{o}{=} \PY{l+s+s2}{\PYZdq{}}\PY{l+s+s2}{Fit: mu = }\PY{l+s+si}{\PYZpc{}.2f}\PY{l+s+s2}{,  std = }\PY{l+s+si}{\PYZpc{}.2f}\PY{l+s+s2}{,  2/sqrt(N) = }\PY{l+s+si}{\PYZpc{}.2f}\PY{l+s+s2}{\PYZdq{}} \PY{o}{\PYZpc{}} \PY{p}{(}\PY{n}{mu}\PY{p}{,} \PY{n}{std}\PY{p}{,} \PY{l+m+mi}{2}\PY{o}{/}\PY{n}{sqrt}\PY{p}{(}\PY{n}{N}\PY{p}{)}\PY{p}{)}
        \PY{n}{plt}\PY{o}{.}\PY{n}{title}\PY{p}{(}\PY{n}{title}\PY{p}{)}
        
        \PY{n}{plt}\PY{o}{.}\PY{n}{xlabel}\PY{p}{(}\PY{l+s+sa}{r}\PY{l+s+s1}{\PYZsq{}}\PY{l+s+s1}{\PYZdl{}}\PY{l+s+s1}{\PYZbs{}}\PY{l+s+s1}{bar}\PY{l+s+si}{\PYZob{}X\PYZcb{}}\PY{l+s+s1}{\PYZdl{}}\PY{l+s+s1}{\PYZsq{}}\PY{p}{)}
        \PY{n}{plt}\PY{o}{.}\PY{n}{ylabel}\PY{p}{(}\PY{l+s+sa}{r}\PY{l+s+s1}{\PYZsq{}}\PY{l+s+s1}{\PYZdl{}Pr(}\PY{l+s+s1}{\PYZbs{}}\PY{l+s+s1}{bar}\PY{l+s+si}{\PYZob{}X\PYZcb{}}\PY{l+s+s1}{)\PYZdl{}}\PY{l+s+s1}{\PYZsq{}}\PY{p}{)}
        \PY{n}{plt}\PY{o}{.}\PY{n}{show}\PY{p}{(}\PY{p}{)}
\end{Verbatim}


    \subsection{Covariância}\label{covariuxe2ncia}


    % Add a bibliography block to the postdoc
    
    
    
    \end{document}
